\chapter{Data Activator (Reflex)}
\index{Data Activator}\index{Reflex}\index{triggers}\index{conditions}\index{actions}\index{alerting!Reflex}\index{alert noise}\index{Teams notifications}%

\begin{objectives}
\begin{itemize}
    \item Explain the Data Activator object model: triggers (now ``rules''), conditions, and actions.
    \item Create Reflex items from Power BI visuals and from Eventstream data.
    \item Configure automated actions: Teams messages, emails, and custom workflows.
    \item Build a sustained-threshold trigger that alerts only after three consecutive minutes of breach.
    \item Suppress alert noise with hysteresis, aggregation, and re-alert policy.
    \item Design a proactive alerting system notifying drivers via custom webhooks when trucks go off route.
\end{itemize}
\end{objectives}

\begin{crosscloud}
Event-driven alerting over analytics---conditions evaluated continuously against live data---exists across clouds, though Fabric's no-code Reflex is unusually approachable.

\vspace{2mm}
{\footnotesize
\begin{tabular}{@{}p{2.8cm}p{3.2cm}p{3.2cm}p{3.2cm}p{3.2cm}@{}}
\toprule
\textbf{Capability} & \textbf{Fabric} & \textbf{AWS} & \textbf{GCP} & \textbf{Snowflake} \\
\midrule
Data-driven alerts & Data Activator (Reflex) & CloudWatch Alarms + SNS & Cloud Monitoring alerts & Snowflake Alerts \\
Stream-native trigger & Reflex over Eventstream/KQL & Kinesis Analytics + Lambda & Dataflow + Eventarc & Streams + tasks + notifications \\
BI-visual trigger & Reflex from Power BI visuals & QuickSight threshold alerts & Looker alerts/schedules & Via BI tool \\
Actions & Teams, email, Power Automate & SNS targets, Lambda & Pub/Sub, Cloud Functions, webhooks & Email, webhook, stored proc \\
Custom workflows & Power Automate flows & Step Functions & Workflows & External functions \\
Evaluation latency & Seconds to minutes & Alarm period granularity & Policy-based & Scheduled evaluation \\
\bottomrule
\end{tabular}}

\vspace{1mm}
Databricks alerts evaluate SQL on schedules; MongoDB Atlas triggers fire on change streams. The portable disciplines are this chapter's heart: \emph{sustained conditions over spikes, actions with owners, and alert noise treated as a reliability bug}.
\end{crosscloud}

\begin{keyterms}
\textbf{Reflex item}: the Fabric item that watches data and acts on conditions. \quad
\textbf{Trigger/rule}: the watched data and the identity of what is being monitored. \quad
\textbf{Condition}: the predicate evaluated as data arrives---thresholds, changes, sustained states. \quad
\textbf{Action}: what fires when the condition holds---Teams, email, Power Automate workflow. \quad
\textbf{Hysteresis}: requiring a condition to hold (or clear) for a sustained period, damping flapping alerts. \quad
\textbf{Alert noise}: the volume of low-value notifications that trains humans to ignore alerts.
\end{keyterms}

\section{Week 15 Roadmap}
Week~15 closes the loop that Chapters~13 and 14 opened: streams are ingested and analyzed; now the platform must \emph{react}. Monday covers the Data Activator object model. Tuesday creates Reflex items from Power BI visuals and Eventstreams. Wednesday covers actions, from Teams messages to custom workflows. Thursday is the hands-on project: an email alert when IoT temperature exceeds a threshold for three consecutive minutes. Friday designs the logistics alerting system with custom app webhooks for off-route trucks.

The week's engineering theme is restraint. Any system can fire an alert; a well-designed system fires the right alert, once, to the right owner, with context---and stays silent otherwise \cite{microsoftDataActivator}.

\begin{notebox}
Alert latency expectation: Reflex evaluates on the order of seconds to minutes depending on source and condition. It is an operational-reaction layer, not a safety-critical control loop. Do not put it between a furnace and a meltdown; do put it between a KPI and the human who can fix it.
\end{notebox}

\section{Monday: Triggers, Conditions, and Actions}
\index{Data Activator!object model}
Data Activator's model has three parts, and every sound design names them explicitly \cite{microsoftDataActivator}:

\begin{definitionbox}
A \textbf{Reflex rule} binds \emph{data} (a stream column or a Power BI measure), a \emph{condition} (a predicate over that data, possibly sustained or aggregated), and an \emph{action} (a notification or workflow fired when the condition holds).
\end{definitionbox}

\begin{itemize}
    \item \textbf{Data and identity}: what is watched (temperature, order count, fraud score) and per-what (per device, per card, per region). The identity granularity defines who can be paged and how many parallel alert streams exist.
    \item \textbf{Conditions}: instant thresholds (temperature $>$ 80), sustained conditions (temperature $>$ 80 for 3 consecutive minutes), and change conditions (value enters/exits a state). Sustained conditions are the primary noise weapon.
    \item \textbf{Actions}: Teams messages to a channel, email to a list, or a Power Automate flow for arbitrary custom workflows---tickets, webhooks, escalations.
\end{itemize}

\begin{pitfall}
Single-spike conditions on noisy telemetry produce alert storms, and alert storms produce ignored alerts, and ignored alerts are worse than no alerts. The ``3 consecutive minutes'' pattern in Thursday's lab is not a detail; it is the difference between a signal and a siren.
\end{pitfall}

\section{Tuesday: Reflex from Power BI and Eventstreams}
\index{Reflex!from Power BI}\index{Reflex!from Eventstream}
Reflex items originate where the data already lives \cite{microsoftDataActivator}:

\subsection{From Power BI Visuals}
A visual on a report---a card showing today's revenue, a gauge showing capacity---can carry a Reflex alert: \emph{set alert} on the visual, define the measure, condition, and cadence, and the rule evaluates as the model refreshes or frames. This path suits KPI-grade reactions: business metrics with business owners.

\subsection{From Eventstreams and KQL}
Over streaming data, Reflex evaluates per event or per aggregation window directly on the stream or the KQL table. This path suits telemetry-grade reactions: device temperatures, error rates, the fraud patterns of the Milestone~4 capstone. The Chapter~14 discipline applies upstream: dedupe before the rule query, because duplicates double-fire alerts.

\begin{tipbox}
Choose the source closest to the truth with the least latency that still reads deduplicated data. For KPIs, that is usually the semantic model; for telemetry, the deduped KQL view---never the raw stream with its replay duplicates.
\end{tipbox}

\section{Wednesday: Actions and Custom Workflows}
\index{Reflex!actions}\index{Power Automate}\index{Teams notifications}
Native actions cover the common cases: Teams channel messages, individual Teams messages, and email. The escape hatch is Power Automate: a Reflex action can trigger a flow, and a flow can do essentially anything---create a ticket in an ITSM tool, call a webhook, post to a third-party chat, open an incident.

Action design is contract design. A good action payload includes: what fired (rule name), on what identity (device, region), the observed values, when, the severity, and the runbook link (Chapter~12's discipline arriving in real time). A payload of ``threshold exceeded'' is a dead end delivered at speed.

Re-alert policy completes the model: after firing, a rule should cool down---suppress repeats while the condition continues to hold, and optionally fire a distinct ``resolved'' notification when it clears. Without cooldown, one sustained breach is fifty identical messages.

\begin{pitfall}
Automating an action that mutates production (restart a service, scale a capacity) from a Reflex rule without a human gate is auto-remediation roulette. Start with notify-only, measure false-positive rate for a month, then graduate specific, reversible actions to automation.
\end{pitfall}

\section{Thursday: Hands-on Project}
\index{hands-on lab}\index{Reflex!hands-on}
The Week~15 lab sets up a Data Activator trigger on the streaming IoT data from Weeks~13--14: an email fires when a device's temperature exceeds a threshold for three consecutive minutes.

\subsection{Goal and Architecture}
Reuse \texttt{es\_iot\_telemetry} and \texttt{kqldb\_iot}. Create a Reflex item \texttt{reflex\_iot\_temp} over the deduplicated KQL view (or the eventstream with a dedup-aware aggregation), with per-device identity, a sustained condition, and an email action.

\subsection{Step-by-Step Actions}
\begin{enumerate}
    \item Create the Reflex item from the eventstream or the KQL table; set the identity to \texttt{DeviceId}.
    \item Define the condition: \texttt{Temperature > 26} sustained for 3 consecutive minutes (per device).
    \item Configure the email action with a complete payload: device, observed values, window, runbook link.
    \item Configure cooldown: suppress repeats for 15 minutes while the condition holds; add a resolved notification.
    \item Inject a synthetic hot streak from the producer (Chapter~14's validation method) for two minutes---no alert should fire.
    \item Extend the streak past three minutes---the alert fires exactly once.
    \item Let the streak end---the resolved notification fires.
\end{enumerate}

\subsection{Validation Checklist}
\begin{itemize}
    \item The rule evaluates per device, not on the fleet aggregate.
    \item A two-minute breach fires nothing; a sustained breach fires once.
    \item The payload contains identity, values, window, severity, and runbook link.
    \item Cooldown suppresses repeats during a continuing breach.
    \item The resolved notification fires on recovery.
    \item The rule reads deduplicated data---a replay drill does not double-fire.
\end{itemize}

\section{Friday: Use Case and Architecture}
\index{alerting!logistics}\index{webhooks!driver notifications}
A logistics company runs 800 delivery trucks with GPS telemetry flowing into Fabric through the Chapter~13 pipeline. Requirement: when a truck deviates from its planned route, notify the driver instantly via the company's custom mobile app---which exposes a webhook---and log the event for the operations team.

\subsection{The Design}
Route deviation is computed in KQL: truck position versus planned route corridor (a stream-to-dimension join against the route table, per Chapter~14). A Reflex rule watches the deduplicated deviation signal per \texttt{TruckId}, with a sustained condition---off-corridor for 60 seconds---to absorb GPS jitter and tunnel dropouts. The action is a Power Automate flow calling the driver-app webhook with a signed payload, plus a Teams message to the regional operations channel. Cooldown prevents nagging; the resolved notification tells the driver they are back on route; every firing writes an audit row to a KQL table for the weekly operations review.

\begin{figure}[H]
    \centering
    \begin{tikzpicture}[node distance=8mm and 9mm]
        \node[doc, minimum width=2.4cm, minimum height=1.0cm] (gps) at (0,0) {Truck GPS\\telemetry};
        \node[stage] (es) at (3.8,0) {Eventstream};
        \node[stage] (kql) at (7.6,0) {KQL: deviation\\(route join)};
        \node[stage, fill=orange!10] (reflex) at (11.6,0) {Reflex rule\\(sustained 60s)};
        \node[doc, minimum width=2.6cm, minimum height=1.0cm] (flow) at (15.4,0) {Power Automate\\$\rightarrow$ app webhook};

        \draw[arrow] (gps) -- (es);
        \draw[arrow] (es) -- (kql);
        \draw[arrow] (kql) -- (reflex);
        \draw[arrow] (reflex) -- (flow);

        \node[note, text width=12.5cm, align=center] at (7.8,-1.9) {Per-truck identity, sustained condition, signed webhook payload, audit row per firing.};
    \end{tikzpicture}
    \caption{Logistics alerting: telemetry to driver notification in seconds.}
    \label{fig:chapter15-logistics}
\end{figure}

\begin{pitfall}
GPS telemetry lies constantly in small ways---jitter, tunnels, urban canyons. A threshold on raw position error alerts on every overpass. Sustained conditions plus a corridor width tuned to measured GPS accuracy are what separate ``driver notified when truly lost'' from ``driver learned to ignore the app.''
\end{pitfall}

\section{Operational Guidance for Week 15}
An alerting estate is a product with users. Measure it: firings per week, false-positive rate, time-to-acknowledge, and the count of alerts nobody acted on. The last metric is the one that predicts disaster.

\subsection{Performance and Reliability}
Keep rule evaluation on deduplicated, materialized data. Bound evaluation windows. Monitor Reflex item health like any production item. Test actions end-to-end monthly---an alert path that silently broke last release is discovered at the worst moment.

\subsection{Security and Governance Checklist}
\begin{itemize}
    \item Webhook payloads signed or token-authenticated; secrets vaulted, never in flow definitions.
    \item Alert destinations (channels, lists) reviewed for information sensitivity---telemetry can leak operational detail.
    \item Every rule has a named owner; ownerless rules are deleted, not orphaned.
    \item Automated mutating actions require documented approval and a rollback path.
    \item Audit rows for every firing kept per the retention policy.
    \item Rules read deduplicated data; a replay must never re-page a human.
\end{itemize}

\section{Interview Q\&A}
\begin{enumerate}
    \item \textbf{Q: Which three parts define a Data Activator trigger?}\\
    A: The data being watched (with its identity granularity), the condition evaluated over it (threshold, sustained, or change), and the action fired when the condition holds.
    \item \textbf{Q: Why require ``3 consecutive minutes'' instead of a single spike?}\\
    A: Sustained conditions filter the transient noise that dominates real telemetry; they convert a flapping siren into a trustworthy signal, which is the only kind humans keep answering.
    \item \textbf{Q: Which actions can Reflex fire natively?}\\
    A: Teams messages and email, plus Power Automate flows as the escape hatch to arbitrary workflows---tickets, webhooks, custom apps.
    \item \textbf{Q: How do you create an alert from a Power BI visual?}\\
    A: Set alert on the visual, binding the rule to the visual's measure; the rule evaluates as the semantic model refreshes or frames, which suits KPI-grade business reactions.
    \item \textbf{Q: What alert latency should you expect from a Reflex item?}\\
    A: Seconds to minutes depending on source and condition cadence---operational reaction time, not hard-real-time control.
    \item \textbf{Q: Why must Reflex rules read deduplicated data?}\\
    A: Because at-least-once delivery and replay produce duplicate events upstream, and a rule over raw duplicates double-fires---paging humans twice for one event.
\end{enumerate}

\section{Further Reading}
Review the Data Activator documentation \cite{microsoftDataActivator}. The upstream disciplines are Chapter~13's eventstream semantics \cite{microsoftFabricEventstreams} and Chapter~14's KQL dedup and windowing \cite{microsoftKQL,microsoftKQLMV}. Alert-as-work-order design inherits Chapter~12's runbook discipline \cite{microsoftFabricMonitoringHub}.

\section*{Key Takeaways}
\begin{itemize}
    \item Data Activator binds data, condition, and action; every sound rule names all three plus an owner.
    \item Sustained conditions are the primary defense against alert storms.
    \item Reflex evaluates in seconds to minutes: operational reaction, not control loops.
    \item Actions carry context---identity, values, severity, runbook---or they are noise delivered at speed.
    \item Cooldown and resolved-notifications complete the alert lifecycle.
    \item Rules read deduplicated data; replay must never re-page a human.
    \item Alert estates are products: measure firings, false positives, and ignored alerts, and delete what nobody acts on.
\end{itemize}

\section{Exercises}
\begin{enumerate}
    \item \textbf{(Conceptual)} Define hysteresis and explain why alerting systems need it, using the truck-corridor example.
    \item \textbf{(Conceptual)} When is a Power-BI-visual rule preferable to a stream-native rule, and vice versa?
    \item \textbf{(Hands-on)} Add a second severity tier to the lab rule (warning at 25, critical at 26) with different actions and cooldowns.
    \item \textbf{(Hands-on)} Build the Power Automate flow posting a signed payload to a test webhook; verify signature validation.
    \item \textbf{(Design)} Design the escalation ladder for the logistics system: driver, regional ops, national ops---with timing and conditions.
    \item \textbf{(Design)} Specify the alert-metrics dashboard: which four metrics, over what window, reviewed how often?
    \item \textbf{(Interview)} In five sentences, explain why ignored alerts are worse than no alerts, and how sustained conditions plus cooldown prevent the slide.
    \item \textbf{(Operational)} Write the monthly alert-path test procedure for the lab's rule chain.
\end{enumerate}

\section{Capstone Project: Fleet Alerting with Noise Discipline}\label{sec:ch15-capstone}
\index{capstone project}\index{Reflex!capstone}\index{alert noise!capstone}

\begin{capstonebox}
\textbf{Mission.} Build the two-tier fleet alerting system: per-device Reflex rules with sustained conditions and cooldowns, a warning/critical severity split, a signed webhook action through Power Automate, a resolved-notification lifecycle, an audit table of every firing, and a one-month noise report proving the estate is worth answering.

\textbf{Estimated time.} 4--6 hours on top of the Weeks 13--14 estate.

\textbf{What you'll practice.}
\begin{itemize}
    \item Sustained-threshold rule design with per-identity granularity.
    \item Severity tiers with differentiated actions.
    \item Signed custom-webhook delivery via Power Automate.
    \item Cooldown, resolution, and audit as first-class lifecycle features.
    \item Measuring alert quality instead of assuming it.
\end{itemize}
\end{capstonebox}

\subsection*{Scenario}
Contoso Devices' support team needs fleet temperature alerting they can trust. The previous system (a cron job emailing on every spike) trained everyone to ignore it within a month. The replacement must fire rarely, fire correctly, carry context, and prove its own quality with data.

\subsection*{Requirements}
\begin{itemize}
    \item Warning and critical rules per device with sustained conditions and distinct actions.
    \item Cooldown and resolved notifications configured and demonstrated.
    \item The Power Automate webhook action with a validated signature.
    \item Audit table capturing every firing: rule, device, values, timestamps, resolution.
    \item Evidence: injected noise produces zero alerts; injected sustained breach produces exactly one per tier.
    \item The noise report template: firings, false positives, acknowledgments, ignored alerts.
\end{itemize}

\subsection*{Implementation Steps}
\begin{enumerate}
    \item Build the warning and critical rules over the deduplicated view.
    \item Configure actions, cooldowns, and resolution notifications.
    \item Build and secure the Power Automate webhook flow.
    \item Implement the audit table write.
    \item Run the noise and breach evidence scenarios; capture the audit rows.
    \item Draft the noise report from the audit table.
\end{enumerate}

\subsection*{Deliverables}
\begin{itemize}
    \item Rule configurations exported or documented with screenshots.
    \item The secured flow definition and a captured signed payload.
    \item Audit table contents across the evidence scenarios.
    \item The noise report template with one populated week (real or simulated).
\end{itemize}

\subsection*{Acceptance Criteria}
\begin{itemize}
    \item[\(\square\)] Injected sub-threshold noise fires nothing.
    \item[\(\square\)] A sustained breach fires exactly once per tier, with full context payloads.
    \item[\(\square\)] Cooldown suppresses repeats; resolution notifies once.
    \item[\(\square\)] Every firing and resolution appears in the audit table.
    \item[\(\square\)] The webhook rejects unsigned payloads in testing.
    \item[\(\square\)] The noise report computes firings, false-positive rate, and ignored-alert count from audit data.
\end{itemize}

\subsection*{Stretch Goals}
\begin{itemize}
    \item Add fleet-level aggregate rules (site temperature, not device) with separate owners.
    \item Add a weekly auto-generated noise report delivered to the rules' owners.
    \item Implement adaptive thresholds from Chapter~14's anomaly detection instead of static values.
    \item Add a suppression window for known maintenance periods, with audit visibility.
\end{itemize}
